\section{Beyond radicals: what can be made constructive}\label{sec:beyond}

The reference pages devote substantial space to Hermite's theta-function
solution, Klein's icosahedral method, hypergeometric differential equations,
and the Doyle--McMullen iteration.  These claims are not excluded by
Abel--Ruffini: they use operations more general than ordinary radicals.  This
section separates their finite local cores from the global analytic work still
required for a complete Lean development.

\subsection{Named hypergeometric functions}

For fixed rational parameters \(a_1,\ldots,a_p,b_1,\ldots,b_q\), define
\[
 c_0=1,
 \qquad
 c_{n+1}=c_n
 \frac{(n+a_1)\cdots(n+a_p)}{(n+b_1)\cdots(n+b_q)(n+1)}.
\]
After clearing the fixed parameter denominators, this is a primitive-recursive
rational recurrence.  The truncated sum
\[
 H_N(z)=\sum_{n<N}c_nz^n
\]
is therefore a finite coded polynomial.  On a named disk where the coefficient
ratio is eventually bounded by a rational \(\theta<1\), the geometric tail
\[
 \left|H(z)-H_N(z)\right|
 \le \frac{|c_Nz^N|}{1-\theta}
\]
gives a primitive-recursive modulus.

\begin{proposition}[Local hypergeometric verification]
For fixed rational parameters, every coefficient identity, linear differential
equation, and algebraic substitution involving a generalized hypergeometric
series can be proved in \ipra{} coefficient by coefficient.  Analytic equality
on a named disk follows when both sides are supplied with effective tail
moduli.
\end{proposition}

\begin{proof}
The coefficient recurrence proves the differential equation after multiplying
by a common denominator and comparing the coefficient of \(z^n\).  Equality of
two named series follows by induction on \(n\) from equal initial coefficients
and a common recurrence.  A geometric or ratio-based tail modulus lifts the
finite equality to the analytic names.
\end{proof}

Theorem~\ref{thm:bring-hypergeom} is a complete example.  The four
hypergeometric basis functions used for all branches of the Bring quintic are
handled the same way: verify the Cockle--Harley recurrence, match four initial
data, and use the zero sum of the five roots.  This is a substantially smaller
task than formalizing a general library of analytic continuation.

The MathWorld observation that equations of the shape
\(x^p+bx^q+c=0\) admit one-variable hypergeometric expansions also has a
finite local interpretation.  Around a simple root or a parameter value where
the equation becomes \(x=X\Phi(x)\), finite Lagrange inversion gives the
coefficient sequence; residue classes of the exponent produce a finite number
of hypergeometric recurrences.  Each fixed pair \((p,q)\) is an \ipra{} theorem
schema.  Uniformity over arbitrary coded \((p,q)\) requires only bounded
factorial and binomial algorithms, not higher-order quantification.

\subsection{Theta series and elliptic data}

For \(|q|<1\), the Jacobi theta constants may be named by
\begin{align*}
 \theta_2(0,q)&=2q^{1/4}\sum_{n\ge0}q^{n(n+1)},\\
 \theta_3(0,q)&=1+2\sum_{n\ge1}q^{n^2}.
\end{align*}
On a named closed disk \(|q|\le r<1\), the omitted tail after \(N\) is bounded
by a geometric series beginning at \(r^{N^2}\).  Thus both theta constants are
primitive-recursively named.  If a rational lower bound for
\(|\theta_3(0,q)|\) is supplied, the modular lambda value
\begin{equation}\label{eq:lambda-theta}
 m(q)=\frac{\theta_2(0,q)^4}{\theta_3(0,q)^4}
\end{equation}
has an effective quotient name.

The complete elliptic integral
\[
 K(k)=\frac\pi2\,{}_2F_1\!\left(\frac12,\frac12;1;k^2\right).
\]
automatically receives a name from the preceding hypergeometric construction
on \(|k|<1\).  The nome
\[
 q=\exp\!\left(-\pi\frac{K(\sqrt{1-k^2})}{K(k)}\right).
\]
gains an effective name when the square roots, denominator lower bound, and
exponential modulus are supplied.  All of these operations are already within
the general effective-analysis architecture used elsewhere in ProveIt.

\begin{theorem}[Finite verification interface for Hermite-type formulas]
\label{thm:hermite-interface}
Fix a compact parameter region avoiding every displayed denominator and branch
cut, and supply:
\begin{enumerate}
\item regular names and moduli for the required theta constants, elliptic
      integrals, roots and exponentials;
\item the finite modular-equation identities relating the auxiliary modulus to
      the Bring parameter;
\item a branch certificate at one base point.
\end{enumerate}
Then the Hermite--Kronecker--Brioschi root expression can be verified in
\ipra{} by approximating both sides and checking the exact algebraic identity
\(x^5-x+a=0\), with uniqueness/continuation on the certified region.
\end{theorem}

\begin{proof}
Every special-function value has a primitive-recursive Cauchy name by the tail
bounds above.  The algebraic combination of names gives a candidate \(x\).
The modular equations and theta identities reduce the residual
\(x^5-x+a\) to zero; at finite precision this is an equality of truncated
series plus an explicit tail bound.  A lower bound for \(|5x^4-1|\) gives local
uniqueness by the quantitative inverse theorem.  Subdivision of a certified
path transports the branch from the base point to the target.
\end{proof}

\begin{frontier}[Global Hermite formula]
A complete proof of the exact formula printed on the source page must formalize
the relevant modular equations, all eighth-root branch choices, and analytic
continuation across a specified cut plane.  The local theta and hypergeometric
objects are fully compatible with \ipra{}, but this archive does not claim that
those global modular identities have already been reduced to Lean certificates.
The status is therefore \(\Outline\), not a rejection of the claim.
\end{frontier}

\subsection{Klein's icosahedral equation}

Klein reduces a principal quintic to an icosahedral invariant equation and
expresses a local inverse by a ratio of two Gauss hypergeometric functions.  The
large polynomial printed by MathWorld is itself an exact finite object.  For a
proposed value of its invariant \(Z\), substitution and polynomial
normalization are immediate \ipra{} tasks.

The hypergeometric ratio has three constructive layers.

\begin{enumerate}
\item Its numerator and denominator are named by rational coefficient
      recurrences, with effective tails for \(|1728Z|<1\).
\item A Wronskian or direct coefficient calculation verifies the associated
      second-order differential equation.
\item The rational icosahedral invariant of the ratio is checked by finite
      series substitution near a base point.  Equality then propagates on a
      certified simply connected region by uniqueness of the differential
      equation.
\end{enumerate}

The local disk is thus accessible with existing tools.  The global theorem that
all relevant branches and icosahedral monodromy are covered requires an
explicit path-and-monodromy library.  Such a library can be entirely
arithmetized: a path is a rational polygon, continuation uses overlapping
rational disks, and a permutation of the five roots is a five-entry table.
What remains is volume, not a need for classical set existence.

\subsection{Differential resolvents in arbitrary fixed degree}

The derivation of Theorem~\ref{thm:differential-resolvent} generalizes.  If
\(P(x,a)=0\) and \(P_x\ne0\), then
\[
 x'=-\frac{P_a}{P_x}.
\]
Repeated total differentiation expresses each \(x^{(j)}\) as a rational
function of \(x,a\).  The vector space over \(\Q(a)\) spanned by
\(1,x,\ldots,x^{d-1}\) has dimension \(d\) modulo a degree-\(d\) equation.
Therefore \(d+1\) derivatives are linearly dependent after reduction modulo
\(P\), producing a linear differential resolvent of order at most \(d\), and
o larger than \(d-1\) when the sum of the branches supplies one dependence.

For fixed \(d\), Gaussian elimination over rational-function codes and
polynomial remainder are primitive recursive.  Thus the source page's claim
that differential resolvents extend to higher-degree algebraic equations has a
uniform finite proof for each fixed degree.  Partial differential resolvents
for several coefficient variables are handled by multi-indexed bounded
derivatives in the same way.

\subsection{Doyle--McMullen iteration}

The source page describes a rational map \(T_Z\) for the Brioschi quintic and
advises iterating from a random start until convergence.  Randomness and the
phrase ``until it converges'' are not a formal theorem.  The correct
constructive interface is a trapping/contraction certificate.

\begin{theorem}[Certified rational iteration]\label{thm:certified-iteration}
Let \(T\) be a rational map with named coefficients.  Suppose a certificate
provides a closed rational disk \(D\), a nonzero lower bound for every
denominator on \(D\), a proof that \(T(D)\subseteq D\), and a rational
\(0\le\lambda<1\) such that
\[
 |T(u)-T(v)|\le\lambda|u-v|
 \quad(u,v\in D).
\]
Then iteration from any named point in \(D\) converges to the unique fixed point
with the primitive-recursive error bound
\[
 |T^n(w)-w_*|
 \le\frac{\lambda^n}{1-\lambda}|T(w)-w|.
\]
If the final algebraic reconstruction formulas are supplied, the resulting
Brioschi roots are effective \ipra{} reals or complexes.
\end{theorem}

\begin{proof}
The usual telescoping proof is a finite induction:
\(|w_{n+1}-w_n|\le\lambda^n|w_1-w_0|\).  A finite geometric sum bounds every
Cauchy tail, and a primitive-recursive logarithmic cutoff makes it smaller than
\(2^{-k}\).  Substitution into the rational map and then the Brioschi
polynomial verifies the limit.
\end{proof}

A global theorem that a generic ``random'' start reaches such a disk is a
separate dynamical statement involving Julia sets and measure.  It is not
needed for certified computation of any fixed instance and is not claimed in
this packet.
